\chapter*{Final Revision Map}
\addcontentsline{toc}{chapter}{Final Revision Map}
\markboth{Final Revision Map}{ }

Use this section after completing all eight chapters. It compresses the course into the comparisons and causal chains most likely to help in revision.

\section*{The complete course in one chain}

\begin{mlfigure}
\begin{tikzpicture}[
  b/.style={draw=MLBlue,fill=MLPaleBlue,rounded corners=2mm,text width=2.45cm,minimum height=1.15cm,align=center,font=\scriptsize\bfseries\color{MLNavy}},
  g/.style={draw=MLGreen,fill=MLPaleTeal,rounded corners=2mm,text width=2.45cm,minimum height=1.15cm,align=center,font=\scriptsize\bfseries\color{MLNavy}},
  o/.style={draw=MLOrange,fill=MLPaleOrange,rounded corners=2mm,text width=2.45cm,minimum height=1.15cm,align=center,font=\scriptsize\bfseries\color{MLNavy}},
  arr/.style={-{Stealth[length=1.8mm]},thick,draw=MLTeal}
]
\node[b] (a) at (-5.2,1.2) {Traditional\\intermediation};
\node[b] (b) at (-2.6,1.2) {Customer\\friction};
\node[g] (c) at (0,1.2) {Digital\\technology};
\node[g] (d) at (2.6,1.2) {Unbundling +\\new platforms};
\node[o] (e) at (5.2,1.2) {Competition +\\new risks};
\node[b] (f) at (-3.9,-1.2) {Financial\\inclusion};
\node[g] (g) at (-1.3,-1.2) {P2P / mobile /\\digital banks};
\node[o] (h) at (1.3,-1.2) {Blockchain +\\crypto};
\node[o] (i) at (3.9,-1.2) {Big Tech +\\ecosystems};
\draw[arr](a)--(b); \draw[arr](b)--(c); \draw[arr](c)--(d); \draw[arr](d)--(e);
\draw[arr](b) to[bend right=24] (f); \draw[arr](c) to[bend right=14] (g); \draw[arr](c) to[bend left=10] (h); \draw[arr](d) to[bend left=24] (i);
\end{tikzpicture}
\end{mlfigure}

\section*{Eight concepts you must distinguish}

\begin{mltable}
\begin{tabularx}{\linewidth}{@{}>{\bfseries\color{MLNavy}}p{0.27\linewidth}X@{}}
\toprule
\rowcolor{MLPaleBlue!92}
Concept & Distinction to remember \\
\midrule
Disintermediation vs unbundling & Disintermediation changes the intermediary; unbundling separates a one-stop provider into specialist services. \\
Unbanked vs underbanked & Unbanked means no formal account; underbanked means partial or insufficient use of mainstream services. \\
Traditional lending vs P2P & Bank lending uses the bank balance sheet and pooling; P2P matches investors and borrowers through a platform. \\
Donation vs reward vs debt vs equity crowdfunding & No return; non-financial reward; repayment with interest; ownership claim. \\
UX vs CX & UX is interface interaction; CX is the entire relationship including reliability, support, security, price, and access. \\
Neobank vs challenger bank & In the Week 5 lecture, no full banking license versus full banking license. \\
Centralized database vs blockchain & Single authoritative state versus replicated, cryptographically linked records governed by network validation/consensus. \\
FinTech vs TechFin & Finance-led technology innovation versus technology-platform-led expansion into financial services. \\
\bottomrule
\end{tabularx}
\end{mltable}

\section*{Recurring FinTech value proposition}

Most cases in the course can be decomposed into five claims:

\begin{enumerate}
  \item \term{Lower cost} through automation, digital distribution, cloud infrastructure, and reduced physical overhead.
  \item \term{Faster service} through fewer manual approval layers and real-time data.
  \item \term{Better customer experience} through intuitive interfaces, transparency, alerts, personalization, and continuous access.
  \item \term{Better matching or risk assessment} through platforms and additional data sources.
  \item \term{Greater access} for customers and businesses who are difficult to serve through legacy economics.
\end{enumerate}

Whenever you make one of these claims in an answer, pair it with a limitation: security, privacy, credit loss, lack of guarantees, regulation, sustainability, integration, or market concentration.

\section*{Recurring technology stack}

\begin{mltable}
\begin{tabularx}{\linewidth}{@{}>{\bfseries\color{MLNavy}}p{0.22\linewidth}X@{}}
\toprule
\rowcolor{MLPaleBlue!92}
Technology & Role in the course \\
\midrule
Mobile & Continuous financial access, notifications, biometrics, camera/OCR, contactless interaction \\
Cloud & Scalable modern infrastructure and faster deployment for digital banks and incumbents \\
APIs & Connect banks to third-party applications and enable platform/open-banking models \\
Big data / analytics & Personalization, credit assessment, fraud detection, and predictive intelligence \\
AI / automation & Faster decisions, conversational assistance, process efficiency, and robo-advice \\
Blockchain & Shared, cryptographically linked transaction history without one authoritative ledger copy \\
Cryptography & Key ownership, transaction authorization, hashing, signatures, and proof-of-work mechanisms \\
\bottomrule
\end{tabularx}
\end{mltable}

\section*{Recurring business-model tension}

The most important business tension across P2P lending, neobanks, Big Tech, and crypto is:

\begin{center}
\term{growth and access} $\quad\longleftrightarrow\quad$ \term{risk, trust, regulation, and sustainable economics}
\end{center}

A good exam answer should explain both sides. Technology can lower one cost while creating another. Removing a bank intermediary can improve price while moving risk to investors. Giving users private-key control increases freedom while reducing recoverability. A neobank can acquire users cheaply while still failing to monetize them sustainably.

\section*{Fast comparison: who owns the customer relationship?}

\begin{mltable}
\begin{tabularx}{\linewidth}{@{}>{\bfseries\color{MLNavy}}p{0.20\linewidth}XXX@{}}
\toprule
\rowcolor{MLPaleBlue!92}
Model & Primary interface & Strategic asset & Main vulnerability \\
\midrule
Traditional bank & Branch + bank-owned digital channels & Trust, license, balance sheet, customer history & Legacy cost and slower change \\
P2P platform & Marketplace & Matching, scoring, low-cost distribution & Credit risk and weak guarantees \\
Digital-only bank & Mobile-first account & CX, modern architecture, data, low servicing cost & Profitability and trust at scale \\
Big Tech / TechFin & Existing super-platform & User base, engagement, data, ecosystem & Regulatory exposure and financial-brand risk \\
Cryptocurrency network & Wallet + protocol & Decentralized transfer and user-controlled keys & Volatility, custody, governance, scalability, regulation \\
\bottomrule
\end{tabularx}
\end{mltable}

\section*{Exam-answer structure}

For an ``analyze'' or ``evaluate'' question, use this order:

\begin{mlsteps}
  \item \term{Define the model precisely.}
  \item \term{Identify the traditional friction or market failure.}
  \item \term{Explain the enabling technology and business-model change.}
  \item \term{State the benefit to each stakeholder.}
  \item \term{State the risk or limitation to each stakeholder.}
  \item \term{Explain the incumbent or regulatory response.}
  \item \term{Conclude with the condition under which the innovation creates sustainable value.}
\end{mlsteps}

\begin{quickcheck}
Without looking back, answer these eight prompts in one or two sentences each:
\begin{enumerate}
  \item Why did the 2008 crisis help create an opening for FinTech?
  \item How do digital financial services change the economics of financial inclusion?
  \item Why can P2P lending serve a borrower rejected by a bank?
  \item Why does a good mobile app still depend on legacy core integration?
  \item Why can a digital-only bank have millions of users and still lose money?
  \item How do hash references make blockchain history tamper-evident?
  \item Why does Big Tech care about financial services even if direct banking profit is modest?
  \item Why is Bitcoin's private-key model both empowering and risky?
\end{enumerate}
\end{quickcheck}

\begin{chapterrecap}
\begin{itemize}
  \item Follow the recurring sequence: friction, digital mechanism, new model, benefit, new risk, institutional response.
  \item The course is not eight isolated topics; each chapter changes either access, intermediation, distribution, infrastructure, or platform control.
  \item Strong answers compare stakeholders and trade-offs rather than listing advantages alone.
  \item Dated market figures in the lecture are evidence for historical context, while the conceptual mechanisms are the main revision target.
\end{itemize}
\end{chapterrecap}
